\documentclass[11pt]{article}

\usepackage[margin=1in]{geometry}
\usepackage{amsmath,amssymb,amsthm,mathtools}
\usepackage{booktabs}
\usepackage{microtype}
\usepackage{enumitem}
\usepackage{hyperref}

\hypersetup{colorlinks=true,linkcolor=blue,citecolor=blue,urlcolor=blue}

\newtheorem{theorem}{Theorem}[section]
\newtheorem{proposition}[theorem]{Proposition}
\newtheorem{lemma}[theorem]{Lemma}
\newtheorem{corollary}[theorem]{Corollary}
\newtheorem{example}[theorem]{Example}
\theoremstyle{definition}
\newtheorem{definition}[theorem]{Definition}
\theoremstyle{remark}
\newtheorem{remark}[theorem]{Remark}

\newcommand{\PP}{\mathcal P}
\newcommand{\Aset}{\mathcal A}
\newcommand{\MM}{\mathcal M}
\newcommand{\Dd}{\mathcal D}
\newcommand{\ZZ}{\mathbb Z}
\newcommand{\CC}{\mathbb C}
\newcommand{\one}{\mathbf 1}
\newcommand{\muM}{\mu}
\newcommand{\ip}[2]{\left\langle #1,#2\right\rangle}
\newcommand{\norm}[1]{\left\lVert #1\right\rVert}
\newcommand{\e}{\mathrm e}

\title{\textbf{Truncated Divisor Bands and the Reciprocal-Emitter Boundary Operator}}
\author{Liang Wang\\
Huazhong University of Science and Technology\\
Wuhan 430074, P.R. China}
\date{August 2026}

\begin{document}
\maketitle

\begin{abstract}
The preceding TPC-211 analysis gives an exact logarithmic M\"obius derivative for a
complete squarefree divisor packet of the literal product-coupled Euler profiles.
The physical transition, however, uses the non-Boolean band $Y_0<d\leq U$ and a
divisor-dependent reciprocal emitter.  This paper isolates that boundary interface.
For any selected divisor family we prove that its common-endpoint leakage is the
signed incidence vector of the selected Boolean faces, with coefficients in the
independent basis $(\log p)_{p\in\PP}$.  The selected packet is exactly the complete
packet minus the missing-subset boundary.  We then define a finite reciprocal
occupancy operator and prove its exact collision Gram identity.  In the natural
direct sum over divisor residue spaces the emitter Gram is block diagonal and full
rank; normalized residuals attain coherent-to-diagonal ratios equal to the number
of divisor blocks.  Thus a divisor cut and reciprocal occupancy, without a theorem
coupling the literal residual profiles across divisors, do not force a saving.  The
finite checks use exact rational arithmetic and a unit-weight reciprocal fixture
($\psi=1$); they are structural certificates, not asymptotic prime-distribution
evidence.  The physical smooth emitter, prime shell, cross-divisor profile theorem,
Gate B, and the twin-prime endpoint remain open.
\end{abstract}

\paragraph{Claim level.}
\texttt{PROVED\_STRUCTURAL\_L1 / STOP\_SCOPED}\\
\texttt{BOUNDARY\_EMITTER}.
No arithmetic $L^2$ credit or fixed-atom credit is claimed.

\section{Introduction}

The working twin-prime route in this repository studies a fixed gap through a
literal transition carrier.  Its central difficulty is not a missing formal
Fourier identity but the simultaneous preservation of the divisor signs, the
physical zero axis, the prime shell, and the cross-divisor Gram structure.  The
local V46 compiler writes the transition residual in the form
\begin{equation}
 \mathfrak R^{\mathrm{AP}}_x
 =-H\sum_{Y_0<d\leq U}\sum_{r\bmod d}
 A_d(r)\widehat{\mathcal R}_d(r),
 \label{eq:physical-residual}
\end{equation}
where, schematically, as in the V46 transition artifact \cite{WangV46},
\begin{equation}
 A_d(r)=\frac{\muM(d)\log d}{d}
 \sum_{q\in\mathcal Q}\sum_{0<|m|\leq dq/H}
 \psi\!\left(\frac{Hm}{dq}\right)
 \mathbf 1_{r\equiv m\overline q\pmod d}.
 \label{eq:emitter}
\end{equation}
The residue vector $\mathcal R_d$ is the literal physical profile at divisor $d$.

TPC-209 showed that a whole-frame Poisson transform leaves divisor-dependent
multiplicative permutations.  TPC-210 showed that independent admissible Poisson
profiles can realize a M\"obius-aligned Gram obstruction \cite{WangTPC210}.  TPC-211 then moved to the
literal product-coupled Euler profiles.  It proved a product cocycle, full divisor
rank, and the exact complete-packet identity \cite{WangTPC211}.
\begin{equation}
 \sum_{\varnothing\ne S\subseteq\PP}
 \muM(d_S)\log(d_S)\Delta_S
 =-\sum_{p\in\PP}\log p\,\mathcal D_p,
 \label{eq:tpc211-derivative}
\end{equation}
where $d_S=\prod_{p\in S}p$ and $\Delta_S$ is the product-profile defect.
The common endpoint cancels on the complete packet when $|\PP|\geq 2$.

The identity in \eqref{eq:tpc211-derivative} does not apply directly to
\eqref{eq:physical-residual}.  The active divisor set is a band in the divisor
lattice, not a complete packet, and $A_d$ changes with $d$.  The precise question
of this paper is therefore:

\begin{quote}
Does the actual divisor cut together with the reciprocal map force a
cross-divisor cancellation before one proves a new theorem for the literal
physical profiles?
\end{quote}

We answer the algebraic question negatively, but in a controlled way.  The answer
has two parts.  First, the cut produces an exact endpoint leakage vector.  Second,
the reciprocal map has an exact collision Gram, but its natural direct-sum geometry
has no cross-divisor entries.  This does not disprove a future estimate on the
coupled physical family.  It identifies the missing input: a coupling theorem must
relate $\mathcal R_d$ across divisors after the boundary and emitter are retained.

Our contributions are:
\begin{enumerate}[leftmargin=2em]
 \item an exact signed-incidence formula for any divisor selector and an exact
       complete-minus-missing decomposition;
 \item an exact reciprocal occupancy collision identity and a block-diagonal
       direct-sum Gram theorem;
 \item a sharp finite alignment construction showing that emitter geometry alone
       cannot supply a universal cross-divisor saving;
 \item exact finite certificates covering four cuts, 5,810 profile coordinates,
       and nine emitter divisor rows.
\end{enumerate}

\paragraph{Scope firewall.}
The proof is finite and algebraic.  We do not estimate the shifted-prime sequence,
the hybrid profile, the smooth weight $\psi$, the prime shell, the four-packet
reassembly, or any asymptotic twin-prime count.  The unit-weight emitter used in
the certificate is explicitly a modeling fixture.

\section{The divisor band as a signed Boolean boundary}

\subsection{Notation}

Let
\[
 \PP=\{p_1,\ldots,p_s\},\qquad p_i>2,
\]
be a finite set of distinct active primes.  For a nonempty subset
$S\subseteq\PP$, put
\[
 d_S=\prod_{p\in S}p,\qquad \varepsilon_S=(-1)^{|S|}.
\]
The squarefree divisors of $t=\prod_{p\in\PP}p$ are indexed by the nonempty
faces of the Boolean cube.  A divisor band is represented by a selector
\[
 \Aset\subseteq 2^{\PP}\setminus\{\varnothing\}.
\]
For the physical interval, the selector is
\begin{equation}
 \Aset_{Y,U}(t)=\{S:Y_0<d_S\leq U\}.
 \label{eq:band-selector}
\end{equation}
The complement within the nonempty Boolean lattice is denoted by
\[
 \MM=(2^{\PP}\setminus\{\varnothing\})\setminus\Aset.
\]

For every active prime define the signed incidence
\begin{equation}
 \eta_p(\Aset)=\sum_{\substack{S\in\Aset\\p\in S}}\varepsilon_S.
 \label{eq:incidence}
\end{equation}
The corresponding logarithmic endpoint coefficient is
\begin{equation}
 L(\Aset)=\sum_{S\in\Aset}\varepsilon_S\log d_S
       =\sum_{p\in\PP}\eta_p(\Aset)\log p.
 \label{eq:log-leakage}
\end{equation}

\begin{lemma}[complete-packet incidence cancellation]
\label{lem:full-incidence}
If $s\geq 2$, then
\[
 \eta_p(2^{\PP}\setminus\{\varnothing\})=0
 \qquad(p\in\PP).
\]
Consequently, the complete packet has $L=0$.
\end{lemma}

\begin{proof}
Fix $p$.  Pair each subset $S$ containing $p$ with the subset obtained by
adding or removing one fixed prime $r\ne p$.  The two signs are opposite.
Equivalently,
\[
 \eta_p=(-1)\sum_{T\subseteq\PP\setminus\{p\}}(-1)^{|T|}
       =-(1-1)^{s-1}=0.
\]
The logarithmic statement follows from \eqref{eq:log-leakage}.
\end{proof}

\begin{theorem}[exact endpoint leakage]
\label{thm:leakage}
For every selector $\Aset$, the common endpoint contribution of the weighted
packet is exactly $L(\Aset)$ times that endpoint.  If $\Aset$ is not complete and
$\eta_p(\Aset)\ne0$ for some $p$, then the endpoint contribution cannot vanish
uniformly over endpoint values.  More precisely, for a scalar endpoint $w$,
\[
 \sum_{S\in\Aset}\varepsilon_S\log(d_S)w=L(\Aset)w.
 \label{eq:endpoint-term}
\]
Moreover, $L(\Aset)=0$ if and only if $\eta_p(\Aset)=0$ for every $p$.
\end{theorem}

\begin{proof}
The first identity is linearity and \eqref{eq:log-leakage}.  Suppose
$L(\Aset)=0$.  Then
\[
 \prod_{p\in\PP}p^{\eta_p(\Aset)}=1
\]
after exponentiating the logarithmic relation.  The exponents are integers,
so unique factorization gives $\eta_p(\Aset)=0$ for all $p$.  The converse is
immediate from \eqref{eq:log-leakage}.  If some incidence is nonzero then the
logarithmic coefficient is nonzero, and choosing $w\ne0$ gives nonzero leakage.
\end{proof}

\begin{example}[the smallest physical band witness]
\label{ex:35-band}
Take $t=5\cdot7=35$ and $Y_0=5$, $U=35$.  Then
\[
 \Aset_{5,35}(35)=\{\{7\},\{5,7\}\}.
\]
The signs are $-1,+1$, so
\[
 \eta_5=1,\qquad \eta_7=0,
 \qquad L(\Aset)= -\log7+\log35=\log5.
 \label{eq:35-leak}
\]
Thus the first proper divisor cut already restores a common endpoint term.
This is an exact witness, not a numerical approximation to a logarithm.
\end{example}

\subsection{Complete minus missing is an operator identity}

Let $\Delta_S$ be any family of vectors in a common real or complex vector
space and let $w$ be a common endpoint vector.  Write
\[
 R_S=w-\Delta_S,
 \qquad
 \mathcal T_\Aset=\sum_{S\in\Aset}\varepsilon_S\log(d_S)R_S.
 \label{eq:packet-functional}
\]
The next statement does not assume the product formula for $\Delta_S$.

\begin{proposition}[boundary decomposition]
\label{prop:boundary-decomp}
With $\MM$ the missing subset family,
\[
 \mathcal T_\Aset=\mathcal T_{\mathrm{full}}-\mathcal T_\MM.
 \label{eq:boundary-decomp}
\]
If the complete packet is a TPC-211 product packet with at least two active
primes, then its endpoint-free part is the marked-prime derivative, while the
missing packet contributes the exact boundary
\[
 \mathcal T_\MM
 =L(\MM)w-\sum_{S\in\MM}\varepsilon_S\log(d_S)\Delta_S.
 \label{eq:missing-boundary}
\]
The endpoint coefficient of the selected packet is $L(\Aset)=-L(\MM)$.
\end{proposition}

\begin{proof}
The first identity is the partition of the complete nonempty subset family
into $\Aset$ and $\MM$.  Expanding $R_S=w-\Delta_S$ gives
\[
 \mathcal T_\MM=L(\MM)w-
 \sum_{S\in\MM}\varepsilon_S\log(d_S)\Delta_S.
\]
Lemma~\ref{lem:full-incidence} gives $L(\Aset)+L(\MM)=0$ when the full packet
has at least two active primes.  No estimate or limiting argument is used.
\end{proof}

\begin{remark}
The boundary in \eqref{eq:missing-boundary} has two distinct pieces: endpoint
leakage and missing profile mass.  Treating the first complete-packet identity
as if it bounded both pieces silently changes the divisor set.  The physical
emitter adds a third issue because its row also changes with $d$.
\end{remark}

\section{The reciprocal occupancy operator}

\subsection{Finite emitter model}

Fix a positive integer $H$, a divisor $d\geq2$, and a finite set $\mathcal Q$
of integers coprime to $d$.  Define the finite index set
\[
 I_d=\left\{(q,m):q\in\mathcal Q,
 0<|m|\leq\left\lfloor\frac{dq}{H}\right\rfloor\right\}.
\]
For unit weights, define the occupancy operator
\[
 E_d:\CC^{I_d}\longrightarrow\CC^{\ZZ/d\ZZ},
 \qquad
 (E_d a)(r)=\sum_{(q,m)\in I_d}
 a(q,m)\mathbf 1_{r\equiv m\overline q\pmod d}.
 \label{eq:occupancy-operator}
\]
The physical emitter is the same pushforward with
\[
 a(q,m)=\psi\!\left(\frac{Hm}{dq}\right),
\]
followed by the nonzero scalar $\muM(d)\log(d)/d$.  We first isolate the
pushforward algebra.

\begin{lemma}[reciprocal collision identity]
\label{lem:collision}
Let $a\equiv1$ on $I_d$ and write $N_d=E_d a$.  Then
\begin{equation}
 \norm{N_d}_2^2
 =\sum_{(q_1,m_1),(q_2,m_2)\in I_d}
 \mathbf 1_{d\mid m_1q_2-m_2q_1}.
 \label{eq:collision}
\end{equation}
More generally, for arbitrary weights $a$,
\begin{equation}
 \norm{E_da}_2^2
 =\sum_{I_d\times I_d}a(q_1,m_1)\overline{a(q_2,m_2)}
 \mathbf 1_{d\mid m_1q_2-m_2q_1}.
 \label{eq:weighted-collision}
\end{equation}
\end{lemma}

\begin{proof}
Expand the squared norm and use
\[
 m_1\overline q_1\equiv m_2\overline q_2\pmod d.
\]
Multiplication by $q_1q_2$, which is invertible modulo $d$, converts this
condition exactly into $m_1q_2\equiv m_2q_1\pmod d$.  The unit-weight case
is the specialization $a=1$.
\end{proof}

\begin{proposition}[direct-sum emitter Gram]
\label{prop:direct-sum}
Let $\mathcal B$ be a finite divisor family and set
\[
 \mathcal H_\mathcal B=\bigoplus_{d\in\mathcal B}\CC^{\ZZ/d\ZZ}.
\]
Let $v_d$ be the unit-weight occupancy vector in the $d$-block, embedded in
this direct sum.  Then
\[
 \ip{v_d}{v_e}=0\quad(d\ne e),
 \qquad
 \ip{v_d}{v_d}=\norm{v_d}_2^2.
 \label{eq:block-gram}
\]
If every $v_d$ is nonzero, the Gram matrix of the family has rank
$|\mathcal B|$.  For any signs $\sigma_d\in\{\pm1\}$ the residuals
\[
 r_d=\sigma_d\frac{v_d}{\norm{v_d}_2^2}
 \label{eq:aligned-residual}
\]
satisfy
\[
 \sigma_d\ip{v_d}{r_d}=1.
 \label{eq:aligned-contribution}
\]
Consequently the coherent-to-diagonal ratio of these unit contributions is
$|\mathcal B|$.
\end{proposition}

\begin{proof}
Different direct-sum blocks are orthogonal, proving \eqref{eq:block-gram}.
Nonzero diagonal entries give full rank.  Equation \eqref{eq:aligned-residual}
gives $\ip{v_d}{r_d}=\sigma_d$, hence \eqref{eq:aligned-contribution}.  The
coherent energy is $|\mathcal B|^2$ and the diagonal energy is
$|\mathcal B|$.
\end{proof}

\begin{remark}
Proposition~\ref{prop:direct-sum} is deliberately a negative interface result.
It does not say that the literal residual family can choose the aligned
residuals independently.  It says that the reciprocal pushforward itself has
not supplied the missing coupling.  A saving theorem must use extra structure
of $\mathcal R_d$, not only the congruence defining $E_d$.
\end{remark}

\section{Exact finite certificates}

\subsection{Protocol}

The certificate in the project directory is generated by exact rational code.
The boundary test uses deterministic rational profiles on the common CRT
coordinate set of size $\prod p$.  This makes the equality in
\eqref{eq:boundary-decomp} a coefficientwise finite identity rather than a
floating-point regression.  The emitter test uses the finite unit-weight
fixture $\psi=1$.  For each divisor the code records the occupancy vector,
its squared norm, and the independent ordered collision count in
\eqref{eq:collision}.

The independent checker reimplements the Boolean and reciprocal arithmetic;
it does not import the producer module.  The optimized Python run is included
as a second execution mode, not as an independent mathematical proof.

\begin{table}[t]
\centering
\caption{Boundary-band certificate cases.  Incidence is ordered by the active
prime list.}
\label{tab:boundary-cases}
\begin{tabular}{@{}c c c c c c@{}}
\toprule
$\PP$ & $Y_0$ & $U$ & active divisors & $\eta(\Aset)$ & coordinates\\
\midrule
$(5,7)$ & 5 & 35 & $7,35$ & $(1,0)$ & 35\\
$(5,7,11)$ & 5 & 35 & $7,35,11$ & $(1,0,-1)$ & 385\\
$(5,7,11)$ & 10 & 77 & $35,11,55,77$ & $(2,2,1)$ & 385\\
$(5,7,11,13)$ & 30 & 100 & $35,55,77,65,91$ & $(3,3,2,2)$ & 5005\\
\bottomrule
\end{tabular}
\end{table}

The first row is the exact witness in Example~\ref{ex:35-band}.  The other
rows show that a divisor band can leave several nonzero incidence coordinates;
the effect is not peculiar to a two-prime packet.  In all four rows the full
incidence vector is zero and the selected plus missing incidence vectors add
to it exactly.

\begin{table}[t]
\centering
\caption{Unit-weight reciprocal-emitter fixtures.  The displayed Gram rank is
the rank in the natural direct sum over the listed divisor blocks.}
\label{tab:emitter-cases}
\begin{tabular}{@{}c c c c c c@{}}
\toprule
divisors & $\mathcal Q$ & $H$ & $\norm{v_d}_2^2$ & rank & ratio\\
\midrule
$(7,35)$ & $(2,3,13)$ & 10 & $(86,470)$ & 2 & 2\\
$(7,11,35,55)$ & $(2,3,13)$ & 10 & $(86,158,470,710)$ & 4 & 4\\
$(35,55,77)$ & $(2,3,13)$ & 10 & $(470,710,1014)$ & 3 & 3\\
\bottomrule
\end{tabular}
\end{table}

The collision counts equal the squared norms entry by entry.  The residuals in
\eqref{eq:aligned-residual}, with the literal M\"obius signs of the displayed
divisors, make all weighted contributions equal to one.  This is a sharp
finite alignment statement for the emitter interface.

\subsection{Certificate status}

The exact certificate covers four boundary cases, 5,810 profile coordinates,
three emitter cases, and nine divisor rows.  The producer, independent checker,
and optimized independent checker all pass.  The following status is part of
the release contract:
\begin{verbatim}
TPC212_ROUTE_ADVANCE = YES
TPC212_STRUCTURAL_THRESHOLD_A = PASS
TPC212_CUT_ENDPOINT_LEAKAGE = PROVED_EXACT
TPC212_BOUNDARY_DECOMPOSITION = PROVED_EXACT
TPC212_RECIPROCAL_COLLISION = PROVED_EXACT_FINITE
TPC212_EMITTER_GRAM = PROVED_EXACT_BLOCK_DIAGONAL
TPC212_EMITTER_ONLY_UNIVERSAL_SAVING = REFUTED_SCOPED
TPC212_LITERAL_PHYSICAL_BOUNDARY_BOUND = OPEN
TPC212_PHYSICAL_CROSS_DIVISOR_GRAM_BOUND = OPEN
TPC212_ARITHMETIC_ADVANCE = NO
TPC212_FIXED_ATOM_CREDIT = 0
TPC212_L2 = NONE
TPC212_FULL_GATE_B_STRICT_1_OVER_400 = UNPAID
\end{verbatim}

\section{What the obstruction does and does not say}

\subsection{The missing physical coupling}

The direct sum in Proposition~\ref{prop:direct-sum} is the correct ambient
space for an arbitrary family of residue vectors $(\mathcal R_d)_{d\in\mathcal
B}$.  It is also exactly the space in which a separate Cauchy estimate sees
one emitter row per divisor.  The proposition shows why that estimate cannot
be upgraded by rhetoric about ``the same reciprocal variable'': after the
residue pushforward, different divisor blocks are still independent unless a
new map identifies them.

For the literal TPC object, such an identifying map might come from the
product-coupled Euler factors, the common endpoint, or the common underlying
integer $u$.  TPC-211 proved that product rank and a shared endpoint do not
provide it automatically.  TPC-212 adds that the cut and reciprocal occupancy
do not provide it either.  The remaining theorem must therefore be of the
form
\begin{equation}
 \left|\sum_{d\in\Aset_{Y,U}(t)}
 \frac{\muM(d)\log d}{d}
 \ip{E_d\psi_d}{\widehat{\mathcal R}_d}\right|
 \leq \text{a saving bound for the coupled family},
 \label{eq:needed-physical-theorem}
\end{equation}
where the right side must use a relation between the different
$\widehat{\mathcal R}_d$.  Applying Cauchy in the direct sum, replacing
$\widehat{\mathcal R}_d$ by arbitrary unit vectors, or replacing the smooth
$\psi_d$ by independent unit weights cannot prove such a relation.

\subsection{Why the finite fixture is not an arithmetic result}

The finite emitter certificate sets $\psi=1$ on a finite set of $(q,m)$ pairs.
It is useful because the congruence collision is then an integer identity and
because every nonzero block can be aligned explicitly.  It is not the physical
smooth emitter, does not include the prime shell, and does not control the
relative phases of the literal residuals.  Therefore the correct label is
\texttt{NUMERICALLY\_CERTIFIED} for the finite rows and
\texttt{MODELING\_CHOICE} for the unit-weight fixture, not an arithmetic
theorem.

The endpoint incidence theorem is stronger in a different direction: it uses
the exact logarithmic basis and unique factorization, so it is a general finite
identity for every cut selector.  Still, it only identifies the term that a
future physical estimate must pay; it does not bound that term in the asymptotic
transition range.

\section{Conclusion and next route}

TPC-212 converts the open phrase ``truncated boundary plus reciprocal emitter''
into two explicit operators.  The divisor cut has a signed Boolean incidence
vector, and the reciprocal map has a collision Gram.  The complete packet
annihilates the first vector, while a proper band generally does not.  The
natural emitter Gram is block diagonal and admits sharp finite alignment.

The strongest positive result is the exact boundary/operator decomposition.  The
strongest obstruction is the scoped failure of emitter-only universal saving.
The next theorem is now sharply specified: construct and bound the coupling map
that takes the literal V46 profile family into these divisor blocks, retaining
the smooth weights and the prime shell before any outer absolute value.  Until
that theorem is proved, Gate B remains open and no twin-prime conclusion follows.

\bibliographystyle{plain}
\bibliography{references}

\end{document}
